% ============================================================
%  LeJEPA Presentation Slides — v2
%  Compile: pdflatex slides_v2.tex && pdflatex slides_v2.tex
% ============================================================
\documentclass[aspectratio=169,11pt]{beamer}

% ============================================================
%  slides_v2_working.tex — work-in-progress
%  Source: slides_v2.tex
%  Applied feedback from: feedback_and_improvement.md
%  Batch progress:
%    [✓] Batch 1 — Slides 1-4
%    [✓] Batch 2 — Slides 5-8
%    [ ] Batch 3 — Slides 9-12
%    [ ] Batch 4 — Slides 13-16
%    [ ] Batch 5 — Verification
% ============================================================

% ---------- Theme ----------
\usetheme{Madrid}
\usecolortheme{default}
\setbeamertemplate{navigation symbols}{}
\setbeamertemplate{footline}[frame number]
\setbeamertemplate{itemize items}[circle]

% ---------- Tiếng Việt có dấu (pdflatex) ----------
\usepackage[utf8]{inputenc}
\usepackage[T5]{fontenc}
\usepackage[vietnamese]{babel}
\usepackage{microtype}

% ---------- Math ----------
\usepackage{amsmath,amssymb,amsthm}
\usepackage{bm}

% ---------- Graphics ----------
\usepackage{tikz}
\usepackage{pgfplots}
\pgfplotsset{compat=1.18}
\usetikzlibrary{
  arrows.meta, positioning, shapes.geometric, shapes.misc,
  fit, backgrounds, decorations.pathreplacing,
  calc, matrix, chains, scopes
}

% ---------- Tables & Boxes ----------
\usepackage{booktabs}
\usepackage{colortbl}
\usepackage{multirow}
\usepackage{array}
\usepackage{tcolorbox}

% ---------- Strikethrough & Symbols ----------
\usepackage[normalem]{ulem}
\usepackage{pifont}
\newcommand{\cmark}{\ding{51}}
\newcommand{\xmark}{\ding{55}}

% ---------- Code ----------
\usepackage{listings}
\lstset{
  basicstyle=\ttfamily\tiny,
  keywordstyle=\color{primary}\bfseries,
  commentstyle=\color{neutral}\itshape,
  backgroundcolor=\color{black!5},
  frame=single, framerule=0pt,
  breaklines=true,
  language=Python,
  extendedchars=true
}

% ---------- Colors — PROFESSIONAL / MUTED (academic style) ----------
\definecolor{primary}{HTML}{2C3E6B}      % Navy — chủ đạo
\definecolor{accent}{HTML}{C8A951}       % Gold nhạt — highlight
\definecolor{bad}{HTML}{B85450}          % Terracotta — limitation
\definecolor{good}{HTML}{5B8C6B}         % Sage — kết quả tốt
\definecolor{neutral}{HTML}{8C8C8C}      % Xám trung
\definecolor{lightblue}{HTML}{E8EAF0}    % nền block chính
\definecolor{lightyellow}{HTML}{F5F0E0}  % nền takeaway
\definecolor{lightgreen}{HTML}{E5EDE8}   % nền exampleblock
\definecolor{lightred}{HTML}{F2E6E5}     % nền alertblock
\definecolor{lightgray}{HTML}{F5F5F5}    % nền trung tính

\setbeamercolor{structure}{fg=primary}
\setbeamercolor{alerted text}{fg=accent}
\setbeamercolor{block title}{bg=primary, fg=white}
\setbeamercolor{block body}{bg=lightblue}
\setbeamercolor{block title alerted}{bg=primary!90, fg=white}  % Navy, KHÔNG đỏ
\setbeamercolor{block body alerted}{bg=lightyellow}
\setbeamercolor{block title example}{bg=good!80, fg=white}
\setbeamercolor{block body example}{bg=lightgreen}

% ---------- Font hierarchy ----------
\setbeamerfont{title}{size=\Large, series=\bfseries}
\setbeamerfont{frametitle}{size=\large, series=\bfseries}
\setbeamerfont{block title}{size=\normalsize, series=\bfseries}

% ---------- Helpers ----------
\newcommand{\takeaway}[1]{%
  \vfill
  \begin{beamercolorbox}[sep=4pt,center,rounded=true]{block title alerted}
    \small\bfseries #1
  \end{beamercolorbox}%
}
\newcommand{\goodtext}[1]{\textcolor{good}{#1}}
\newcommand{\badtext}[1]{\textcolor{bad}{#1}}
\newcommand{\hl}[1]{\textcolor{accent}{\bfseries #1}}

% ---------- TikZ styles dùng chung toàn deck ----------
\tikzset{
  % Boxes cơ bản
  stdbox/.style={
    draw, rounded corners=4pt,
    minimum height=0.8cm, minimum width=1.6cm,
    font=\small, inner sep=5pt
  },
  % Màu theo ngữ nghĩa — muted, chuyên nghiệp
  primarybox/.style={stdbox, fill=primary!85, text=white, draw=primary},
  accentbox/.style={stdbox, fill=accent!40,  text=primary, draw=accent!70},
  goodbox/.style={stdbox,    fill=lightgreen, draw=good!60},
  badbox/.style={stdbox,     fill=lightred,   draw=bad!60},
  neutralbox/.style={stdbox, fill=lightgray,  draw=neutral!50},
  lightbluebox/.style={stdbox, fill=lightblue, draw=primary!30},
  % Mũi tên
  mainarrow/.style={->, thick, primary},
  subarrow/.style={->, thin,  neutral},
  dashedarrow/.style={->, dashed, thin, neutral},
  % Node tròn nhỏ (embedding)
  embnode/.style={circle, draw=good!60, fill=good!15, minimum size=0.6cm, font=\small}
}

% ============================================================
\title[LeJEPA]{\textbf{LeJEPA}: Provable \& Scalable\\
Self-Supervised Learning Without the Heuristics}
\author{Randall Balestriero \& Yann LeCun (2025)}
\institute{arXiv: 2511.08544}
\date{}
% ============================================================

\begin{document}

% ============================================================
%  SLIDE 1 — Title (Layout A)
% ============================================================
\begin{frame}[plain]
\vspace{0.5cm}
\begin{columns}[c]
\column{0.44\textwidth}
  \begin{center}
    {\Huge\bfseries\color{primary} LeJEPA}\\[6pt]
    {\large\itshape Provable \& Scalable SSL — Without the Heuristics}\\[14pt]
    {\small Randall Balestriero \& Yann LeCun (2025)}\\[4pt]
    {\small\color{neutral} arXiv: 2511.08544}\\[4pt]
    \vspace*{12pt}
    {\footnotesize Presenter: [CHỦ-SLIDE-ĐIỀN]}
  \end{center}
\column{0.54\textwidth}
  \centering
  \begin{tikzpicture}[every node/.style={font=\tiny}]
    \node[lightbluebox, minimum width=1.0cm, minimum height=0.6cm] (img) at (0,0)   {Ảnh};
    \node[primarybox,   minimum width=1.0cm, minimum height=0.6cm] (enc) at (2.0,0) {$f_\theta$};
    \node[embnode, minimum size=0.55cm] (em1) at (3.8, 0.0) {};
    \node[embnode, minimum size=0.45cm] (em2) at (4.3, 0.7) {};
    \node[embnode, minimum size=0.50cm] (em3) at (4.4,-0.6) {};
    \node[embnode, minimum size=0.40cm] (em4) at (5.0, 0.3) {};
    \node[embnode, minimum size=0.45cm] (em5) at (4.9,-1.0) {};
    \draw[dashed, color=good!50, thick] (4.5, -0.2) circle (1.4cm);
    \node[font=\tiny, color=good] at (5.75, 1.05) {$\mathcal{N}(0,I)$};
    \node[font=\tiny, below=4pt of img] {Dữ liệu};
    \node[font=\tiny, below=4pt of enc] {Bộ mã hóa};
    \draw[mainarrow] (img) -- (enc);
    \draw[->, thick, color=good] (enc) -- (em1);
    \draw[->, thick, color=good] (enc) -- (em2);
    \draw[->, thick, color=good] (enc) -- (em3);
    \draw[->, thick, color=good] (enc) -- (em4);
    \draw[->, thick, color=good] (enc) -- (em5);
  \end{tikzpicture}
\end{columns}
\takeaway{LeJEPA: chứng minh lý thuyết được, triển khai trong 50 dòng Python}
\end{frame}

% ============================================================
%  SLIDE 2 — Ba paradigm học máy (3-col, Layout C mở rộng)
% ============================================================
\begin{frame}{Ba Paradigm Học Máy}
\begin{columns}[t]
\column{0.31\textwidth}
  \begin{block}{Supervised}
    \centering\small Ảnh $\to$ Nhãn\\[4pt]
    \footnotesize
    \textcolor{bad}{\xmark}\ Cần annotation\\
    \textcolor{bad}{\xmark}\ Đắt, khan hiếm\\
    \textcolor{bad}{\xmark}\ Không scale
  \end{block}
\column{0.31\textwidth}
  \begin{block}{Reconstruction}
    \centering\small Ảnh $\to$ Pixel\\[4pt]
    \footnotesize
    \textcolor{bad}{\xmark}\ Lãng phí vào pixel\\
    \textcolor{bad}{\xmark}\ Không học semantic\\
    \textcolor{bad}{\xmark}\ Compute cao
  \end{block}
\column{0.31\textwidth}
  \begin{exampleblock}{Self-Supervised}
    \centering\small View1 $\leftrightarrow$ View2\\[4pt]
    \footnotesize
    \textcolor{good}{\cmark}\ Tự sinh tín hiệu\\
    \textcolor{good}{\cmark}\ Không cần nhãn\\
    \textcolor{good}{\cmark}\ Scale tốt\\[3pt]
    \textbf{\color{good} ← LeJEPA ở đây}
  \end{exampleblock}
\end{columns}
\takeaway{Annotation đắt — cấu trúc dữ liệu thì miễn phí}
\end{frame}

% ============================================================
%  SLIDE 3 — Collapse (Layout C)
% ============================================================
\begin{frame}{Collapse Là Lý Do SSL Thường Thất Bại}
\begin{columns}[t]
\column{0.47\textwidth}
  \begin{center}
    \textcolor{bad}{\xmark\ \textbf{Nghiệm tầm thường}}
  \end{center}
  \centering
  \begin{tikzpicture}[every node/.style={font=\tiny}]
    % Vùng embedding — nền đỏ gợi "không gian bị thu về 1 điểm"
    \draw[dashed, color=neutral!40, fill=lightred, opacity=0.7] (2.3,-0.1) circle (1.3cm);
    \node[font=\tiny, color=neutral, align=center] at (2.3, 1.25) {không gian $\mathbb{R}^K$};
    % Inputs
    \node[stdbox, fill=lightred, draw=bad, minimum width=0.8cm, minimum height=0.5cm] (x1) at (0, 1.0) {$x_1$};
    \node[stdbox, fill=lightred, draw=bad, minimum width=0.8cm, minimum height=0.5cm] (x2) at (0, 0.3) {$x_2$};
    \node[stdbox, fill=lightred, draw=bad, minimum width=0.8cm, minimum height=0.5cm] (x3) at (0,-0.4) {$x_3$};
    \node[stdbox, fill=lightred, draw=bad, minimum width=0.8cm, minimum height=0.5cm] (x4) at (0,-1.1) {$x_4$};
    % Caption PHÍA TRÊN collapse point
    \node[font=\small\bfseries, color=bad, align=center] at (2.3, 0.85)
      {$f(x_1)\!=\!\cdots\!=\!c$};
    % Collapse point — to, đậm
    \node[circle, fill=bad, text=white, minimum size=1.0cm, font=\bfseries]
      (pt) at (2.3,-0.1) {$\mathbf{c}$};
    % Curved arrows
    \draw[->, thick, color=bad!70] (x1.east) to[bend left=10]  (pt.west);
    \draw[->, thick, color=bad!70] (x2.east) to[bend left=5]   (pt.west);
    \draw[->, thick, color=bad!70] (x3.east) to[bend right=5]  (pt.west);
    \draw[->, thick, color=bad!70] (x4.east) to[bend right=10] (pt.west);
  \end{tikzpicture}
\column{0.47\textwidth}
  \begin{center}
    \textcolor{good}{\cmark\ \textbf{Mục tiêu thực}}
  \end{center}
  \centering
  \begin{tikzpicture}[every node/.style={font=\tiny}]
    % Vùng embedding — nền xanh gợi "phủ đều không gian"
    \draw[dashed, color=good!50, fill=lightgreen, opacity=0.7] (2.3,-0.1) circle (1.3cm);
    \node[font=\tiny, color=neutral, align=center] at (2.3, 1.25) {không gian $\mathbb{R}^K$};
    % Inputs
    \node[stdbox, fill=lightgreen, draw=good, minimum width=0.8cm, minimum height=0.5cm] (y1) at (0, 1.0) {$x_1$};
    \node[stdbox, fill=lightgreen, draw=good, minimum width=0.8cm, minimum height=0.5cm] (y2) at (0, 0.3) {$x_2$};
    \node[stdbox, fill=lightgreen, draw=good, minimum width=0.8cm, minimum height=0.5cm] (y3) at (0,-0.4) {$x_3$};
    \node[stdbox, fill=lightgreen, draw=good, minimum width=0.8cm, minimum height=0.5cm] (y4) at (0,-1.1) {$x_4$};
    % Scatter 2D: z_i phân bố trong vùng embedding
    \node[embnode, minimum size=0.5cm] (p1) at (1.6,  0.7) {$z_1$};
    \node[embnode, minimum size=0.5cm] (p2) at (2.9,  0.5) {$z_2$};
    \node[embnode, minimum size=0.5cm] (p3) at (1.7, -0.8) {$z_3$};
    \node[embnode, minimum size=0.5cm] (p4) at (2.8, -0.6) {$z_4$};
    % Thin arrows (1-1 mapping nhưng tỏa ra 2D)
    \draw[->, thin, color=good!60] (y1.east) -- (p1.west);
    \draw[->, thin, color=good!60] (y2.east) -- (p2.west);
    \draw[->, thin, color=good!60] (y3.east) -- (p3.west);
    \draw[->, thin, color=good!60] (y4.east) -- (p4.west);
  \end{tikzpicture}
\end{columns}
\takeaway{Cần ràng buộc chống collapse — nhưng ràng buộc NÀO là đúng?}
\end{frame}

% ============================================================
%  SLIDE 4 — JEPA (Layout B: 60% TikZ trái, 40% text phải)
% ============================================================
\begin{frame}{JEPA Học Trạng Thái Trừu Tượng — Không Phải Pixel}
\begin{columns}[c]
\column{0.62\textwidth}
  \centering
  \begin{tikzpicture}[every node/.style={font=\tiny}, scale=0.82, transform shape]
    % Ảnh gốc
    \node[lightbluebox, minimum width=0.9cm, minimum height=0.55cm] (src) at (0.0, 0.0) {$x$};
    \node[font=\tiny, below=3pt of src] {Ảnh gốc};
    % Views
    \node[lightbluebox, minimum width=0.9cm, minimum height=0.55cm] (v1) at (1.4,  0.7) {view 1};
    \node[lightbluebox, minimum width=0.9cm, minimum height=0.55cm] (v2) at (1.4, -0.7) {view 2};
    \draw[subarrow] (src.north east) -- (v1.west);
    \draw[subarrow] (src.south east) -- (v2.west);
    % Encoders
    \node[primarybox, minimum width=0.9cm, minimum height=0.55cm] (e1) at (2.8,  0.7) {$f_\theta$};
    \node[primarybox, minimum width=0.9cm, minimum height=0.55cm] (e2) at (2.8, -0.7) {$f_\theta$};
    \draw[mainarrow] (v1) -- (e1);
    \draw[mainarrow] (v2) -- (e2);
    % Latents
    \node[accentbox, minimum width=0.7cm, minimum height=0.55cm] (z1) at (4.0,  0.7) {$z_1$};
    \node[accentbox, minimum width=0.7cm, minimum height=0.55cm] (z2) at (4.0, -0.7) {$z_2$};
    \draw[mainarrow] (e1) -- (z1);
    \draw[mainarrow] (e2) -- (z2);
    % Predictor + ẑ_1 — giãn đều, đủ thở
    \node[primarybox, minimum width=1.1cm, minimum height=0.55cm] (pred) at (5.4, -0.7) {Dự báo};
    \node[accentbox,  minimum width=0.7cm, minimum height=0.55cm] (zh)   at (6.8, -0.7) {$\hat{z}_1$};
    \draw[mainarrow] (z2)   -- (pred);
    \draw[mainarrow] (pred) -- (zh);
    % Loss — cuối hàng dưới
    \node[goodbox, fill=good!60, text=white, font=\small\bfseries,
          minimum width=1.1cm, minimum height=0.65cm] (loss) at (8.2, -0.7) {Loss};
    \draw[mainarrow] (zh) -- (loss);
    % z1 → loss.north: ngang hàng trên, thả thẳng xuống
    \draw[mainarrow] (z1.east) -- (8.2, 0.7) -- (loss.north);
    % Shared encoder background
    \begin{pgfonlayer}{background}
      \node[draw=neutral, dashed, fill=neutral!5, rounded corners=4pt,
            fit=(e1)(e2), inner sep=5pt] (encbox) {};
    \end{pgfonlayer}
    \node[font=\tiny, color=neutral, above=3pt of encbox] {chia sẻ trọng số};
  \end{tikzpicture}
\column{0.36\textwidth}
  \begin{block}{\footnotesize 2 Nguyên Tắc}
    \footnotesize
    \textbf{1.\ Dự báo được:}\\
    $f_\theta(x_{t+1})$ dự báo từ $f_\theta(x_t)$\\[4pt]
    \textbf{2.\ Không collapse:}\\
    Embedding phải đa dạng
  \end{block}
  \vspace{4pt}
  \begin{alertblock}{\footnotesize JEPA vs Pixel Models}
    \footnotesize
    \textcolor{bad}{\xmark}\ Pixel models (MAE, PixelCNN):\\
    \quad dự báo pixel $\to$ lãng phí vào nhiễu\\[3pt]
    \textcolor{good}{\cmark}\ JEPA: dự báo trạng thái\\
    \quad $\to$ học cấu trúc ngữ nghĩa
  \end{alertblock}
\end{columns}
\takeaway{JEPA học "thế giới thay đổi thế nào", không học "pixel trông như thế nào"}
\end{frame}

% ============================================================
%  SLIDE 5 — Whac-A-Mole (Layout F: bảng 5 method × 3 cột)
% ============================================================
\begin{frame}{Heuristic Stack: Vá Lỗ, Không Có Lý Thuyết}
\centering\footnotesize
\renewcommand{\arraystretch}{1.3}
\begin{tabular}{l >{\raggedright\arraybackslash}p{2.8cm} c >{\raggedright\arraybackslash}p{3.0cm}}
\toprule
\textbf{Phương pháp} & \textbf{Cơ chế chống collapse} & \textbf{\#HP} & \textbf{Hạn chế} \\
\midrule
SimCLR  & Negative samples (contrastive)          & 3   & \cellcolor{bad!20}\textcolor{bad}{\xmark}\ $\mathcal{O}(N^2)$ bộ nhớ \\
BYOL    & Kiến trúc bất đối xứng + stop-grad      & 5   & \cellcolor{bad!20}\textcolor{bad}{\xmark}\ Không có lý thuyết \\
DINO    & Teacher--student + EMA schedule          & 7+  & \cellcolor{bad!20}\textcolor{bad}{\xmark}\ $>$5 siêu tham số \\
VICReg  & Làm trắng phương sai (variance)          & 4   & \cellcolor{bad!20}\textcolor{bad}{\xmark}\ Under-specified \\
I-JEPA  & Stop-gradient + mask prediction          & 5   & \cellcolor{bad!20}\textcolor{bad}{\xmark}\ Nhạy cảm, chỉ ViT \\
\midrule
\rowcolor{good!15}
\textbf{LeJEPA} & \textbf{Iso Gaussian} & \textbf{1} & \textcolor{good}{\cmark}\;Provable \\
\bottomrule
\end{tabular}
\takeaway{Không có framework nào từ first principles — chỉ là empirical patches}
\end{frame}


% ============================================================
%  SLIDE 6 — Câu hỏi pivot (Layout B: arrow diagram phải, text trái)
% ============================================================
\begin{frame}{Câu Hỏi Thay Đổi Mọi Thứ}
\begin{columns}[c]
\column{0.38\textwidth}
  \begin{block}{Cách tiếp cận cũ}
    \footnotesize
    \textcolor{bad}{\xmark}\ ``Làm sao \textbf{tránh} collapse?''
  \end{block}
  \vspace{10pt}
  \begin{exampleblock}{Câu hỏi LeJEPA}
    \footnotesize
    \textcolor{good}{\cmark}\ ``Distribution \textbf{nào} là tối ưu?''
  \end{exampleblock}
\column{0.60\textwidth}
  \centering
  \begin{tikzpicture}[every node/.style={font=\small}]
    % 3 node chính theo chiều dọc
    \node[badbox, minimum width=3.2cm, minimum height=1.0cm,
          text width=3.0cm, align=center]
      (old) at (0, 1.5) {\footnotesize Prevent\\Collapse};
    \node[primarybox, minimum width=3.2cm, minimum height=1.0cm,
          text width=3.0cm, align=center]
      (pivot) at (0, 0.0) {\footnotesize\bfseries Distribution\\Tối Ưu?};
    \node[goodbox, minimum width=3.2cm, minimum height=1.0cm,
          text width=3.0cm, align=center]
      (result) at (0,-1.5) {\footnotesize $\mathcal{N}(0,I)$\\Optimal};
    % Mũi tên
    \draw[->, thick, color=bad!70, dashed] (old) -- (pivot);
    \draw[->, thick, color=good]
      (pivot) -- node[right=5pt, font=\tiny, color=neutral] {Theorem 3.3} (result);
  \end{tikzpicture}
\end{columns}
\takeaway{Đổi câu hỏi: ``tránh collapse'' $\to$ ``đạt distribution tối ưu'' — giải từ gốc}
\end{frame}

% ============================================================
%  SLIDE 7 — Linear probe setup (Layout D: equation center + 2 case)
% ============================================================
\begin{frame}{Framework: Đánh Giá Qua Linear Probe Ridge Regression}
\begin{center}
  \[
  \hat{\beta}
  = \underset{\beta}{\arg\min}\;
    \underbrace{\|\mathbf{y} - \mathbf{Z}\beta\|_2^2}_{\text{prediction error}}
  + \underbrace{\lambda\|\beta\|_2^2}_{\text{regularization}}
  \]
\end{center}
\vspace{0.15cm}
\begin{columns}[t]
\column{0.48\textwidth}
  \centering
  \textcolor{bad}{\textbf{\large \xmark\ Anisotropic}}\\[6pt]
  \parbox[b][1.4cm][c]{\linewidth}{\centering
    \begin{tikzpicture}[scale=0.55, every node/.style={font=\tiny}]
      \draw[bad, thick, fill=bad!10] (0,0) ellipse (1.8cm and 0.55cm);
      \foreach \px/\py in {1.5/0.1, -1.4/-0.1, 0.9/0.42, -0.8/-0.40, 1.1/-0.32}
        {\fill[bad!70] (\px,\py) circle (2.5pt);}
      \node[bad, font=\tiny, below=6pt] at (0,-0.55) {$\lambda_1 \gg \lambda_2 \gg \cdots \gg \lambda_K$};
    \end{tikzpicture}
  }\\[4pt]
  \footnotesize
  \textcolor{bad}{$\Rightarrow$ Neighborhood bị méo}\\[2pt]
  \textcolor{bad}{$\Rightarrow$ Estimator lean vào chiều lớn}\\[6pt]
  \textbf{\color{bad}\large Bias \& Variance $\uparrow$}
\column{0.48\textwidth}
  \centering
  \textcolor{good}{\textbf{\large \cmark\ Isotropic}}\\[6pt]
  \parbox[b][1.4cm][c]{\linewidth}{\centering
    \begin{tikzpicture}[scale=0.55, every node/.style={font=\tiny}]
      \draw[good, thick, fill=good!10] (0,0) circle (0.9cm);
      \foreach \px/\py in {0.7/0.58, -0.8/0.42, 0.85/-0.28, -0.55/-0.72, 0.18/0.88}
        {\fill[good!80] (\px,\py) circle (2.5pt);}
      \node[good, font=\tiny, below=6pt] at (0,-0.90) {$\lambda_1 = \lambda_2 = \cdots = \lambda_K$};
    \end{tikzpicture}
  }\\[4pt]
  \footnotesize
  \textcolor{good}{$\Rightarrow$ Neighborhood cân bằng}\\[2pt]
  \textcolor{good}{$\Rightarrow$ Estimator unbiased}\\[6pt]
  \textbf{\color{good}\large Bias \& Variance $\downarrow$}
\end{columns}
\vspace{0.35cm}
\takeaway{Cùng ``energy'' (tổng phương sai) — hình học khác nhau $\Rightarrow$ hiệu quả khác nhau}
\end{frame}

% ============================================================
%  SLIDE 8 — Lemma aniso (Layout C: 2 panel pgfplots + 2 alertblock)
% ============================================================
\begin{frame}{Bằng Chứng 1: Aniso Tệ Hơn Iso Về Cả Bias Và Variance}
\begin{columns}[t]
% --- Panel trái: Bias ---
\column{0.48\textwidth}
  \begin{tikzpicture}
  \begin{axis}[
    width=5.2cm, height=3.0cm,
    title={Bias theo số mẫu $N$},
    title style={font=\footnotesize\bfseries},
    xlabel={Số mẫu $N$},
    ylabel={Bias ($\uparrow$ = worse)},
    xlabel style={font=\tiny},
    ylabel style={font=\tiny, yshift=-2pt},
    tick label style={font=\tiny},
    xmin=50, xmax=500,
    ymin=0.0, ymax=0.70,
    legend pos=outer north east,
    legend style={font=\tiny, draw=none, fill=lightgray, inner sep=2pt},
    legend cell align={left},
    grid=major, grid style={gray!20},
  ]
    % Aniso: bias cao hơn, giảm chậm — nằm TRÊN iso → khớp Bias(aniso) > Bias(iso)
    \addplot[color=bad, thick, mark=none, smooth]
      coordinates {(50,0.62)(100,0.48)(200,0.33)(300,0.24)(400,0.17)(500,0.13)};
    \addlegendentry{Aniso \xmark}
    % Iso: bias thấp hơn, giảm nhanh — nằm DƯỚI aniso
    \addplot[color=good, thick, mark=none, smooth]
      coordinates {(50,0.39)(100,0.26)(200,0.15)(300,0.10)(400,0.06)(500,0.04)};
    \addlegendentry{Iso \cmark}
  \end{axis}
  \end{tikzpicture}
  \begin{center}\scriptsize
    \textbf{Lemma 3.1:} Với $\lambda>0$, $\exists\,\mathbf{y}$:
    $\text{Bias}(\hat\beta_{\text{aniso}}) > \text{Bias}(\hat\beta_{\text{iso}})$
  \end{center}

% --- Panel phải: Variance ---
\column{0.48\textwidth}
  \begin{tikzpicture}
  \begin{axis}[
    width=5.2cm, height=3.0cm,
    title={Variance của decision boundary},
    title style={font=\footnotesize\bfseries},
    xlabel={Dim.~1},
    ylabel={Dim.~2},
    xlabel style={font=\tiny},
    ylabel style={font=\tiny, yshift=-2pt},
    tick label style={font=\tiny},
    xmin=-2.8, xmax=2.8,
    ymin=-2.2, ymax=2.2,
    legend pos=outer north east,
    legend style={font=\tiny, draw=none, fill=lightgray, inner sep=2pt},
    legend cell align={left},
    grid=major, grid style={gray!20},
  ]
    % Aniso — plot đầu lấy legend, còn lại forget plot
    \addplot[color=bad, thin, domain=-2.8:2.8, samples=2] {-0.40*x + 0.50};
    \addlegendentry{Aniso \xmark}
    \addplot[color=bad, thin, forget plot, domain=-2.8:2.8, samples=2] { 0.20*x - 0.30};
    \addplot[color=bad, thin, forget plot, domain=-2.8:2.8, samples=2] { 0.55*x + 0.10};
    \addplot[color=bad, thin, forget plot, domain=-2.8:2.8, samples=2] { 0.80*x - 0.60};
    \addplot[color=bad, thin, forget plot, domain=-2.8:2.8, samples=2] { 0.00*x + 0.70};
    % Iso — plot đầu lấy legend, còn lại forget plot
    \addplot[color=good, thick, domain=-2.8:2.8, samples=2] {-0.55*x + 0.05};
    \addlegendentry{Iso \cmark}
    \addplot[color=good, thin, forget plot, domain=-2.8:2.8, samples=2] {-0.50*x + 0.10};
    \addplot[color=good, thin, forget plot, domain=-2.8:2.8, samples=2] {-0.57*x - 0.05};
    \addplot[color=good, thin, forget plot, domain=-2.8:2.8, samples=2] {-0.52*x + 0.00};
  \end{axis}
  \end{tikzpicture}
  \begin{center}\scriptsize
    \textbf{Lemma 3.2:}
    $\text{tr}(\text{Var}(\hat\beta_{\text{aniso}})) > \text{tr}(\text{Var}(\hat\beta_{\text{iso}}))$
  \end{center}
\end{columns}
\takeaway{Iso thắng aniso về CẢ bias lẫn variance — với mọi downstream task}
\end{frame}

% ============================================================
%  SLIDE 9 — Theorem iso optimal (Layout E: alertblock + intuition + visual)
% ============================================================
\begin{frame}{Bằng Chứng 2: Iso Gaussian Tối Ưu Cho kNN Và Kernel Regression}
\begin{alertblock}{Theorem 3.3 (Balestriero \& LeCun, 2025)}
  \tiny
  Với $\text{Tr}(\text{Cov}(\mathbf{Z})) = \kappa$, iso Gaussian là \textbf{unique minimizer}:
  \[
  \text{ISB}_{k\text{-NN}}
  = \underbrace{\frac{r_0^4}{(K+2)^2}\tau_g^2 J(p)}_{\text{bias do hình học}}
  + O(r_0^4),\quad
  J(p) = \underbrace{\text{Fisher info}(p)}_{\min\text{ tại }\mathcal{N}(0,I)}
  \]
  $\Rightarrow$ $\arg\min_p \text{ISB}(p) = \mathcal{N}(\mathbf{0},\tfrac{\kappa}{K}I)$ \textbf{(unique)}
\end{alertblock}
\begin{columns}[c]
\column{0.56\textwidth}
  \footnotesize
  \textbf{Intuition — tại sao $\mathcal{N}(0,I)$?}\\[2pt]
  kNN và kernel regression dựa vào \textbf{khoảng cách Euclidean}.\\[4pt]
  \textcolor{bad}{\xmark}\ Aniso: neighborhood lệch — prediction bị bias\\[2pt]
  \textcolor{good}{\cmark}\ Iso: neighborhood cân bằng — unbiased\\[4pt]
  {\tiny\color{neutral} Fisher info $J(p)$ minimize duy nhất tại Gaussian.}
\column{0.42\textwidth}
  \centering
  \begin{tikzpicture}[every node/.style={font=\tiny}, scale=0.78]
    \node[font=\tiny\bfseries, color=bad] at (1.3, 2.0)
      {\xmark\ Aniso: bị lệch};

    \draw[color=bad!40, fill=bad!8, thick]
      (0.6,1.4) ellipse (1.0cm and 0.40cm);
    \fill[bad] (0.6,1.4) circle (2pt)
      node[above=2pt] {query};
    \fill[neutral!60] (1.3,1.42) circle (1.5pt);
    \fill[neutral!60] (1.5,1.38) circle (1.5pt);
    \node[font=\tiny\bfseries, color=good] at (1.3, 0.5)
      {\cmark\ Iso: cân bằng};
    \draw[color=good!50, fill=good!8, thick]
      (1.1,0.0) circle (0.58cm);
    \fill[good] (1.1,0.0) circle (2pt)
      node[above=2pt] {query};
    \fill[neutral!60] (0.65, 0.22) circle (1.5pt);
    \fill[neutral!60] (1.55, 0.22) circle (1.5pt);
    \fill[neutral!60] (1.10,-0.40) circle (1.5pt);
  \end{tikzpicture}
\end{columns}
\takeaway{Kết quả UNIQUE: không distribution nào khác đạt được ISB tối thiểu này}
\end{frame}

% ============================================================
%  SLIDE 10 — Geometric intuition (Layout C: ellipse aniso vs sphere iso)
% ============================================================
\begin{frame}{Geometric Intuition: Hình Cầu Là Tối Ưu}
\centering
\begin{tikzpicture}[every node/.style={font=\small}]
  % ---- Panel trái: Aniso (ellipse) ----
  \begin{scope}[xshift=-3.8cm]
    \node[color=bad, font=\footnotesize\bfseries] at (0, 2.0)
      {\xmark\ Anisotropic};
    \draw[color=bad, fill=bad!8, thick]
      (0,0) ellipse (1.6cm and 0.6cm);
    \draw[->, thick, color=bad]   (0,0) -- (1.6, 0);
    \draw[->, thick, color=bad!50](0,0) -- (0,  0.6);
    \node[font=\tiny, right=2pt] at (1.6,0) {$\lambda_1$ lớn};
    \node[font=\tiny, above=2pt] at (0, 0.6) {$\lambda_2$};
    \foreach \px/\py in
      {1.4/0.1, 1.1/-0.25, -1.3/0.2, -1.0/-0.28,
       0.5/0.48, -0.4/0.52} {
      \fill[bad!50] (\px,\py) circle (1.8pt);
    }
    \node[color=bad, font=\tiny, align=center] at (0,-1.1)
      {Khoảng cách méo\\prediction bị bias};
  \end{scope}
  % ---- Ký hiệu ở giữa ----
  \node[font=\Large, color=neutral] at (0, 0) {$\Rightarrow$};
  % ---- Panel phải: Iso (tròn) ----
  \begin{scope}[xshift=3.8cm]
    \node[color=good, font=\footnotesize\bfseries] at (0, 2.0)
      {\cmark\ Isotropic};
    \draw[color=good, fill=good!8, thick]
      (0,0) circle (1.0cm);
    \draw[->, thick, color=good] (0,0) -- (1.0,  0.0);
    \draw[->, thick, color=good] (0,0) -- (0.0,  1.0);
    \draw[->, thick, color=good] (0,0) -- (-0.71,-0.71);
    \node[font=\tiny, right=2pt] at (1.0,0) {$\lambda_1$};
    \node[font=\tiny, above=2pt] at (0, 1.0) {$\lambda_2$};
    \node[font=\tiny, below left=1pt] at (-0.71,-0.71) {$\lambda_k$};
    \foreach \px/\py in
      {0.9/0.35, 0.45/0.9, -0.8/0.62, -1.0/0.1,
       -0.62/-0.78, 0.82/-0.54} {
      \fill[good!70] (\px,\py) circle (1.8pt);
    }
    \node[color=good, font=\tiny, align=center] at (0,-1.3)
      {Khoảng cách đều\\prediction unbiased};
  \end{scope}
\end{tikzpicture}

\takeaway{$\mathcal{N}(0,I)$: ``bản đồ phẳng'' \xmark\ Aniso=méo \cmark\ Iso=đều — unbiased mọi hướng}
\end{frame}




% ============================================================
%  SLIDE 11 — Tóm tắt lý thuyết (Layout B: flow diagram trái + bảng phải)
% ============================================================
\begin{frame}{Kết Luận Lý Thuyết: Design Principle Cho LeJEPA}
\begin{columns}[c]
% --- Cột trái: flow diagram ---
\column{0.50\textwidth}
  \centering
  \begin{tikzpicture}[every node/.style={font=\tiny}]
    % Node đầu: bất kỳ downstream task
    \node[lightbluebox, minimum width=3.2cm, minimum height=0.65cm,
          text width=3.0cm, align=center]
      (task) at (0, 1.5) {Bất kỳ downstream task};
    % Node giữa: 3 probe types
    \node[neutralbox, minimum width=1.7cm, minimum height=0.55cm]
      (lp) at (-1.9, 0.5) {Linear probe};
    \node[neutralbox, minimum width=1.7cm, minimum height=0.55cm]
      (knn) at (0.0, 0.5) {kNN probe};
    \node[neutralbox, minimum width=1.7cm, minimum height=0.55cm]
      (kp) at (1.9, 0.5) {Kernel probe};
    % Node kết quả
    \node[goodbox, minimum width=3.2cm, minimum height=0.65cm,
          text width=3.0cm, align=center]
      (opt) at (0,-0.5) {Optimal khi $\mathbf{Z}\sim\mathcal{N}(0,I)$};
    % Node design principle
    \node[primarybox, minimum width=3.4cm, minimum height=0.65cm,
          text width=3.2cm, align=center]
      (dp) at (0,-1.5) {Design: $f_\theta(x)\sim\mathcal{N}(0,I)$};
    % Arrows
    \draw[mainarrow] (task.south) -- (lp.north);
    \draw[mainarrow] (task.south) -- (knn.north);
    \draw[mainarrow] (task.south) -- (kp.north);
    \draw[mainarrow] (lp.south)  -- (opt.north);
    \draw[mainarrow] (knn.south) -- (opt.north);
    \draw[mainarrow] (kp.south)  -- (opt.north);
    \draw[mainarrow, very thick] (opt.south) -- (dp.north);
    % Label mũi tên cuối
    \node[right=3pt of opt.south, font=\tiny, color=neutral, xshift=2pt, yshift=-8pt] {Theorem 3.3};
  \end{tikzpicture}

% --- Cột phải: bảng so sánh ---
\column{0.48\textwidth}
  \begin{block}{So sánh với Prior SSL}
    \footnotesize
    \renewcommand{\arraystretch}{1.05}
    \begin{tabular}{>{\raggedright\arraybackslash}p{2.1cm}
                    >{\raggedright\arraybackslash}p{1.6cm}
                    >{\raggedright\arraybackslash}p{1.6cm}}
    \toprule
    \textbf{Tiêu chí}
      & \textbf{Prior SSL}
      & \textbf{LeJEPA} \\
    \midrule
    Lý do chọn distribution
      & \cellcolor{bad!15}\tiny ``It works''
      & \cellcolor{good!15}\tiny\bfseries Proven optimal \\
    Guarantee
      & \cellcolor{bad!15}\tiny Không có
      & \cellcolor{good!15}\tiny\bfseries Unique min. \\
    Linear probe
      & \cellcolor{bad!15}\tiny Chưa rõ
      & \cellcolor{good!15}\tiny\bfseries \cmark \\
    Nonlinear probe
      & \cellcolor{bad!15}\tiny Chưa rõ
      & \cellcolor{good!15}\tiny\bfseries \cmark \\
    \bottomrule
    \end{tabular}
  \end{block}
  \begin{exampleblock}{\footnotesize Ý nghĩa}
    \footnotesize
    Lần đầu tiên SSL có \textbf{lý do lý thuyết} cho việc chọn distribution — không phải thử nghiệm.
  \end{exampleblock}
\end{columns}
\takeaway{Lần đầu tiên: biết CHÍNH XÁC embeddings nên phân phối thế nào}
\end{frame}

% ============================================================
%  SLIDE 12 — Thách thức high-dim (Layout B)
% ============================================================
\begin{frame}{Thách Thức: Kiểm Tra Distribution Trong Không Gian Cao Chiều}
\begin{columns}[c]
\column{0.55\textwidth}
  \centering
  \begin{tikzpicture}[x=1.3cm, y=0.7cm, every node/.style={font=\tiny}]
    \draw[->, neutral] (0,0) -- (3.5,0) node[right] {Complexity};
    \draw[neutral] (0,0) -- (0,-4.5);
    
    % MMD
    \node[left, align=right] at (0,-1) {MMD};
    \fill[bad!80] (0,-1.2) rectangle (3,-0.8);
    \node[right, bad] at (3,-1) {$\mathcal{O}(N^2)$};
    
    % KL Divergence
    \node[left, align=right] at (0,-2) {KL Divergence};
    \fill[bad!80] (0,-2.2) rectangle (1.5,-1.8);
    \node[right, bad] at (1.5,-2) {Unstable};
    
    % Direct Multivariate
    \node[left, align=right] at (0,-3) {Direct Multivariate};
    \fill[bad!80] (0,-3.2) rectangle (3,-2.8);
    \node[right, bad] at (3,-3) {$\mathcal{O}(N^2)$};
    
    % SIGReg
    \node[left, align=right] at (0,-4) {\textbf{SIGReg}};
    \fill[good!80] (0,-4.2) rectangle (0.8,-3.8);
    \node[right, good] at (0.8,-4) {\textbf{$\mathcal{O}(N)$}};
  \end{tikzpicture}
\column{0.43\textwidth}
  \begin{block}{3 Yêu Cầu Bắt Buộc}
    \footnotesize
    \begin{itemize}
      \item[\textcolor{good}{\cmark}] \textbf{Differentiable:}\\
        Cho gradient-based training
      \item[\textcolor{good}{\cmark}] \textbf{$\mathcal{O}(N)$ complexity:} \\
        Scale lên millions samples
      \item[\textcolor{good}{\cmark}] \textbf{Provably correct:}\\
        Guarantee convergence đến $\mathcal{N}(0,I)$
    \end{itemize}
  \end{block}
  \vspace{4pt}
  \begin{alertblock}{Vấn Đề}
    \footnotesize Mọi multivariate normality test chuẩn (BHEP, Mardia) đều $\mathcal{O}(N^2)$.
  \end{alertblock}
\end{columns}
\takeaway{Không thể test $\mathcal{N}(0,I)$ trực tiếp trong high-dim — cần sketching}
\end{frame}

% ============================================================
%  SLIDE 13 — Random Slicing (Layout B)
% ============================================================
\begin{frame}{Giải Pháp: Sketch Qua Random 1D Projections}
\begin{columns}[c]
\column{0.55\textwidth}
  \centering
  \begin{tikzpicture}[scale=0.9, every node/.style={font=\tiny}]
    % 3D Cloud
    \draw[primary!30, fill=primary!10, thick] (0, 0) ellipse (1.2cm and 0.8cm);
    \node[font=\footnotesize, primary] at (0, 1.0) {Cloud Điểm $\mathbb{R}^K$};
    \foreach \px/\py in {-0.5/0.2, 0.3/0.4, 0.1/-0.3, -0.6/-0.1, 0.5/-0.2} {
      \fill[primary] (\px,\py) circle (1.5pt);
    }
    
    % Arrows
    \draw[->, thick, accent] (1.0, 0.2) -- (2.5, 1.2) node[midway, above left] {$\mathbf{u}_1$};
    \draw[->, thick, accent] (1.2, 0) -- (2.5, 0) node[midway, above] {$\mathbf{u}_2$};
    \draw[->, thick, accent] (0.8,-0.4) -- (2.5,-1.2) node[midway, below left] {$\mathbf{u}_M$};
    
    % 1D Histograms
    \begin{scope}[xshift=3.0cm, yshift=1.2cm, scale=0.4]
      \draw[neutral] (-1.5,0) -- (1.5,0);
      \draw[good, fill=good!20, thick] (-1,0) sin (0,1) cos (1,0);
      \node[right, good] at (1.2,0.5) {Khớp $\mathcal{N}(0,1)$};
    \end{scope}
    
    \begin{scope}[xshift=3.0cm, yshift=0cm, scale=0.4]
      \draw[neutral] (-1.5,0) -- (1.5,0);
      \draw[bad, fill=bad!20, thick] (-1,0) sin (-0.5,0.8) cos (0,0) sin (0.5, 0.5) cos (1,0);
      \node[right, bad] at (1.2,0.5) {Không khớp};
    \end{scope}
    
    \begin{scope}[xshift=3.0cm, yshift=-1.2cm, scale=0.4]
      \draw[neutral] (-1.5,0) -- (1.5,0);
      \draw[good, fill=good!20, thick] (-1,0) sin (0,1) cos (1,0);
      \node[right, good] at (1.2,0.5) {Khớp $\mathcal{N}(0,1)$};
    \end{scope}
  \end{tikzpicture}
\column{0.43\textwidth}
  \begin{alertblock}{Hyperspherical Cramér-Wold}
    \footnotesize
    $X \stackrel{d}{=} Y \iff \langle \mathbf{u}, X \rangle \stackrel{d}{=} \langle \mathbf{u}, Y \rangle, \quad \forall \mathbf{u} \in \mathbb{S}^{K-1}$
  \end{alertblock}
  \vspace{5pt}
  \begin{block}{Phiên bản thực tế}
    \footnotesize
    Không cần \textbf{tất cả} các hướng — chỉ cần lấy mẫu $M$ hướng ngẫu nhiên và tái sử dụng qua nhiều bước huấn luyện.
  \end{block}
  \vspace{5pt}
  \footnotesize
  \textbf{Theorem 4.2:}\\
  Với $M$ hướng ngẫu nhiên, test vẫn consistent nếu $M \to \infty$ cùng training.
\end{columns}
\takeaway{Bài toán K-chiều $\to$ $M$ bài toán 1-chiều độc lập}
\end{frame}

% ============================================================
%  SLIDE 14 — Moment Family (Layout C modified)
% ============================================================
\begin{frame}{Test 1D Họ Moment: Mạnh Nhưng Không Stable}
\begin{columns}[t]
\column{0.48\textwidth}
  \begin{block}{Extended Jarque-Bera (EJB)}
    \footnotesize
    \[
    \resizebox{0.95\linewidth}{!}{$
    \text{EJB}(\mathbf{u}) = \frac{N\hat\mu^2}{\hat\sigma^2} 
    + \frac{(N-1)(\hat\sigma^2 - 1)^2}{2} 
    + \frac{N}{6}\left(\text{skew}^2 + \frac{(\text{kurt}-3)^2}{4}\right)
    $}
    \]
  \end{block}
  \vspace{5pt}
  \begin{alertblock}{Theorem 4.3}
    \footnotesize
    Minimize $\sum c_k(m_k(P) - m_k(Q))^2$ với $K$ hữu hạn \textbf{không} imply $P = Q$.
  \end{alertblock}
  \footnotesize
  \begin{itemize}
      \item \textbf{$K$ nhỏ:} shortcut solutions tồn tại.
      \item \textbf{$K$ lớn:} gradient explode.
  \end{itemize}
\column{0.48\textwidth}
  \begin{center}
    \footnotesize\textbf{Gradient Explosion Visualization}
  \end{center}
  \centering
  \begin{tikzpicture}
    \begin{axis}[
      width=5.5cm, height=3.5cm,
      xlabel={Moment order $k$},
      ylabel={Gradient norm $\|\partial L / \partial z_i\|$},
      xlabel style={font=\tiny},
      ylabel style={font=\tiny, yshift=-2pt},
      tick label style={font=\tiny},
      xmin=1, xmax=6,
      ymin=0, ymax=100,
      grid=major, grid style={gray!20},
    ]
      \addplot[color=bad, thick, domain=1:6, samples=100] {2^(x)};
      \fill[bad!20, opacity=0.5] (axis cs:4,0) rectangle (axis cs:6,100);
      \node[bad, font=\tiny, align=center] at (axis cs:5, 50) {Unstable\\gradient\\region};
    \end{axis}
  \end{tikzpicture}
\end{columns}
\takeaway{Moment-based: identifiability vs stability — không thể có cả hai}
\end{frame}

% ============================================================
%  SLIDE 15 — CDF Family (Layout C)
% ============================================================
\begin{frame}{Test 1D Họ CDF: Chính Xác Nhưng Không Differentiable}
\begin{columns}[t]
\column{0.50\textwidth}
  \begin{block}{Công Thức CDF Test}
    \footnotesize
    \[
    T_w = N\int_{-\infty}^{\infty}(F_N(x) - F(x))^2 w(x)\, dF(x)
    \]
    \begin{itemize}
        \item $w(x) = 1 \implies$ Cramér-von Mises
        \item $w(x) = [F(x)(1-F(x))]^{-1} \implies$ Anderson-Darling
    \end{itemize}
  \end{block}
  \begin{alertblock}{2 Hạn Chế Chính}
    \scriptsize
    \textbf{1.\ Sort breaks parallelism:}\\
    Multi-GPU cần synchronize $\to$ bottleneck.\\
    \textbf{2.\ Non-differentiable:}\\
    Cần relaxation (differentiable sort) $\to$ thêm hyperparameter.
  \end{alertblock}
\column{0.46\textwidth}
  \begin{center}
    \footnotesize\textbf{Sort Operation Visualization}
  \end{center}
  \centering
  \begin{tikzpicture}[scale=0.7, transform shape, every node/.style={font=\tiny}]
    % Unsorted array
    \node[stdbox, fill=lightgray] (u1) at (0, 2) {4.2};
    \node[stdbox, fill=lightgray] (u2) at (0, 1.4) {1.1};
    \node[stdbox, fill=lightgray] (u3) at (0, 0.8) {3.5};
    \node[stdbox, fill=lightgray] (u4) at (0, 0.2) {2.8};
    
    % Sort arrow
    \draw[->, thick, primary] (1, 1.1) -- (2, 1.1) node[midway, above] {Sort};
    
    % Sorted array
    \node[stdbox, fill=lightgreen] (s1) at (3, 2) {1.1};
    \node[stdbox, fill=lightgreen] (s2) at (3, 1.4) {2.8};
    \node[stdbox, fill=lightgreen] (s3) at (3, 0.8) {3.5};
    \node[stdbox, fill=lightgreen] (s4) at (3, 0.2) {4.2};
    
    % Text
    \node[bad, font=\footnotesize, align=center] at (1.5, -0.6) {$\mathcal{O}(N \log N)$, non-parallel};
    \node[bad, font=\footnotesize, align=center] at (1.5, -1.2) {\underline{$\partial(\text{sort})/\partial x = \text{undefined}$}};
  \end{tikzpicture}
\end{columns}
\takeaway{CDF-based: cần sort $\to$ không parallel + non-differentiable $\to$ loại}
\end{frame}

% ============================================================
%  SLIDE 16 — Epps-Pulley (Layout D)
% ============================================================
\begin{frame}{Epps-Pulley: Stable, Scalable, Provable}
\begin{center}
  \footnotesize
  \[
  EP = N \int_{-\infty}^{\infty} \underbrace{|\hat\phi_X(t) - \phi(t)|^2}_{\text{ECF vs target CF}} 
  \underbrace{w(t)}_{\text{Gaussian window}} dt
  \]
  với $\hat\phi_X(t) = \frac{1}{N}\sum_{j=1}^N e^{itX_j}, \quad \phi(t) = e^{-t^2/2} \text{ (\textbf{CF của } } \mathcal{N}(0,1)\text{)}$
\end{center}
\vspace{5pt}
\begin{columns}[T]
\column{0.48\textwidth}
  \vspace{0pt}
  \centering
  \begin{tikzpicture}
    \begin{axis}[
      width=5.5cm, height=3.2cm,
      xlabel={$t$},
      ylabel={CF Value},
      xlabel style={font=\tiny},
      ylabel style={font=\tiny, yshift=-2pt},
      tick label style={font=\tiny},
      xmin=-5, xmax=5,
      ymin=-0.5, ymax=1.2,
      grid=major, grid style={gray!20},
    ]
      \addplot[color=primary, thick, domain=-5:5, samples=100] {exp(-x^2/2)};
      \addplot[color=good, thin, domain=-5:5, samples=100] {exp(-x^2/2) + 0.1*sin(deg(5*x))*exp(-x^2/4)};
      \node[font=\tiny, primary] at (axis cs:2.5, 0.8) {Target CF};
      \node[font=\tiny, good] at (axis cs:-2.5, 0.8) {Empirical CF};
    \end{axis}
  \end{tikzpicture}
  \begin{itemize}
      \item[\textcolor{good}{\cmark}] \footnotesize \textbf{Differentiable} (ECF là trung bình của $e^{itx}$)
      \item[\textcolor{good}{\cmark}] \footnotesize \textbf{$\mathcal{O}(N)$} tính toán (tổng hợp, không cần sort)
      \item[\textcolor{good}{\cmark}] \footnotesize \textbf{DDP-friendly} (chỉ cần \texttt{all\_reduce})
  \end{itemize}
\column{0.48\textwidth}
  \vspace{0pt}
  \begin{alertblock}{Bounded Gradient (Theorem 4.5)}
    \footnotesize
    \[
    \resizebox{0.9\linewidth}{!}{$
    \left|\frac{\partial EP}{\partial z_i}\right| \leq \frac{4\sigma^2}{N}, \qquad
    \left|\frac{\partial^2 EP}{\partial z_i^2}\right| \leq \frac{C\sqrt{\pi}\sigma^3}{2N}
    $}
    \]
    Gradient và curvature đều \textbf{BỊ CHẶN} — bất kể phân phối của dữ liệu.
  \end{alertblock}
\end{columns}
\takeaway{ECF bị chặn trong $[-1,1] \implies$ gradient không bao giờ explode}
\end{frame}

% ============================================================
%  SLIDE 17 — SIGReg formula (Layout B)
% ============================================================
\begin{frame}[fragile]{SIGReg: Sketched Isotropic Gaussian Regularization}
\begin{columns}[c]
\column{0.50\textwidth}
  \begin{block}{Định Nghĩa (SIGReg)}
    \footnotesize
    \[
    \resizebox{0.95\linewidth}{!}{$
    \text{SIGReg} = \frac{1}{|\mathcal{A}|}\sum_{\mathbf{a}\in\mathcal{A}} 
    \underbrace{EP\big(\{\mathbf{a}^\top f_\theta(\mathbf{x}_n)\}_{n=1}^N\big)}_{\text{Epps-Pulley trên 1D projection}}
    $}
    \]
  \end{block}
  \vspace{5pt}
  \centering
  \begin{tikzpicture}[every node/.style={font=\tiny}, scale=0.85]
    \node[goodbox, minimum width=2.0cm, align=center] (p1) at (0,0) {Differentiable\\(ECF = avg)};
    \node[goodbox, minimum width=2.0cm, align=center] (p2) at (2.8,0) {$\mathcal{O}(N)$ Scalable\\(no sort)};
    \node[goodbox, minimum width=2.0cm, align=center] (p3) at (5.6,0) {Provably Correct\\(Cramér-Wold)};
    \draw[thick, neutral] (p1) -- (p2);
    \draw[thick, neutral] (p2) -- (p3);
  \end{tikzpicture}
\column{0.48\textwidth}
\begin{lstlisting}
def SIGReg(x, global_step, num_slices=1024):
    g = torch.Generator(device=x.device)
    g.manual_seed(global_step) # sync seed across GPUs
    A = torch.randn(x.size(1), num_slices, generator=g)
    A /= A.norm(p=2, dim=0) # unit vectors
    t = torch.linspace(0, 3, 17) # integration points
    phi = torch.exp(-0.5 * t**2) # N(0,1) CF
    x_t = (x @ A).unsqueeze(-1) * t
    cos_m = x_t.cos().mean(0) # empirical CF (real)
    sin_m = x_t.sin().mean(0) # empirical CF (imag)
    err = (cos_m - phi).square() + sin_m.square()
    return (err @ weights).mean() * N
\end{lstlisting}
\end{columns}
\takeaway{$\sim$50 dòng code PyTorch, chạy trên bất kỳ GPU/TPU nào}
\end{frame}

% ============================================================
%  SLIDE 18 — Curse of dimensionality (Layout C)
% ============================================================
\begin{frame}{SIGReg Thoát Khỏi Curse of Dimensionality}
\begin{columns}[t]
\column{0.48\textwidth}
  \begin{center}
    \footnotesize\textbf{Fixed vs Resampled Directions}
  \end{center}
  \centering
  \begin{tikzpicture}
    \begin{axis}[
      width=5.5cm, height=3.5cm,
      xlabel={Số hướng $M$},
      ylabel={Test Error},
      xlabel style={font=\tiny},
      ylabel style={font=\tiny, yshift=-2pt},
      tick label style={font=\tiny},
      xmin=0, xmax=1024,
      ymin=0, ymax=1,
      legend style={font=\tiny, at={(0.95,0.95)}, anchor=north east},
      grid=major, grid style={gray!20},
    ]
      \addplot[color=bad, thick, domain=16:1024, samples=50] {0.8 * exp(-x/500) + 0.1};
      \addlegendentry{Fixed}
      \addplot[color=good, thick, domain=16:1024, samples=50] {0.8 * exp(-x/100) + 0.1};
      \addlegendentry{Resampled}
    \end{axis}
  \end{tikzpicture}
  \vspace{-5pt}
  \begin{block}{\footnotesize Giải Pháp 1: SGD Resampling}
    \footnotesize Qua $T$ bước huấn luyện, tổng số hướng = $M \times T \implies$ coverage tích lũy.
  \end{block}
\column{0.48\textwidth}
  \begin{center}
    \footnotesize\textbf{Error Bound vs Smoothness (Theorem 4.7)}
  \end{center}
  \centering
  \begin{tikzpicture}
    \begin{axis}[
      width=5.5cm, height=3.5cm,
      xlabel={Sobolev smoothness $\alpha$},
      ylabel={Required $M$ cho error $\varepsilon$},
      xlabel style={font=\tiny},
      ylabel style={font=\tiny, yshift=-2pt},
      tick label style={font=\tiny},
      xmin=1, xmax=5,
      ymin=0, ymax=1000,
      grid=major, grid style={gray!20},
    ]
      \addplot[color=primary, thick, domain=1:5, samples=50] {1000 * x^(-2)};
    \end{axis}
  \end{tikzpicture}
  \vspace{-5pt}
  \begin{block}{\footnotesize Giải Pháp 2: Tính Trơn (Smoothness)}
    \footnotesize DNN naturally smooth (Batch Norm, Dropout) $\implies \alpha$ lớn $\implies M = \mathcal{O}(K)$ là đủ.
  \end{block}
\end{columns}
\takeaway{Nhờ tính trơn + resampling: $M=16$ đủ dù $K=2048$}
\end{frame}

% ============================================================
%  SLIDE 19 — LeJEPA Loss (Layout D + G)
% ============================================================
\begin{frame}{LeJEPA = Prediction Loss + SIGReg}
\vspace{-5pt}
\begin{center}
  \tiny
  \[
  \mathcal{L}_\text{LeJEPA} = 
  \underbrace{(1-\textcolor{accent}{\lambda})\cdot \frac{1}{B}\sum_{n=1}^B \mathcal{L}_\text{pred}(\{\mathbf{z}_{n,v}\})}_{\textcolor{primary}{\text{Invariance: các views phải agree}}} 
  + \underbrace{\frac{\textcolor{accent}{\lambda}}{V}\sum_{v=1}^V \text{SIGReg}(\{\mathbf{z}_{n,v}\})}_{\textcolor{good}{\text{SIGReg: push về } \mathcal{N}(\mathbf{0},I)}}
  \]
  \vspace{-10pt}
  \[
  \mathcal{L}_\text{pred} = \frac{1}{V}\sum_{v'=1}^V \|\boldsymbol{\mu}_n - \mathbf{z}_{n,v'}\|_2^2, \quad \boldsymbol{\mu}_n = \frac{1}{V_g}\sum_{v=1}^{V_g}\mathbf{z}_{n,v}
  \]
\end{center}
\vspace{-5pt}
\centering
\begin{tikzpicture}[
    scale=0.45, transform shape,
    box/.style={rectangle, draw, rounded corners, align=center, minimum height=0.6cm, font=\scriptsize},
    every node/.style={font=\scriptsize}
  ]
  % Inputs
  \node (img) at (0, 3) {Image $x$};
  
  % Views
  \node[box] (v1) at (3, 4) {view 1\\(global)};
  \node[box] (v2) at (3, 2.5) {view 2\\(global)};
  \node (vd) at (3, 1.5) {$\vdots$};
  \node[box] (v8) at (3, 0) {view 8\\(local)};
  
  \draw[->, neutral] (img) |- (v1) node[pos=0.75, above, font=\tiny] {aug};
  \draw[->, neutral] (img) |- (v2) node[pos=0.75, above, font=\tiny] {aug};
  \draw[->, neutral] (img) |- (v8) node[pos=0.75, above, font=\tiny] {aug};
  
  % Encoder
  \node[box, fill=neutral!10, minimum height=4.5cm] (enc) at (6, 2) {Encoder $f_\theta$\\ (bất kỳ)};
  \draw[->, neutral] (v1) -- (v1 -| enc.west);
  \draw[->, neutral] (v2) -- (v2 -| enc.west);
  \draw[->, neutral] (v8) -- (v8 -| enc.west);
  
  % Projector
  \node[box, fill=neutral!10, minimum height=4.5cm] (proj) at (8.5, 2) {Proj Head\\ (3-layer MLP)};
  \draw[->, neutral] (enc) -- (proj);
  
  \node (z) at (10.5, 2) {$\mathbf{z}$};
  \draw[->, neutral] (proj) -- (z);
  
  % Losses
  \node[box, draw=primary, fill=primary!10, text=primary] (inv) at (7, -1.5) {Invariance Loss\\$(\mathbf{z}_v \approx \boldsymbol{\mu})$};
  \node[box, draw=good, fill=good!10, text=good] (sig) at (11, -1.5) {SIGReg Loss\\$(\mathbf{z} \sim \mathcal{N}(0, I))$};
  
  \draw[->, primary, thick] (z.south) -- (10.5, -0.5) -| (inv.north);
  \draw[->, good, thick] (z.south) -- (10.5, -0.5) -| (sig.north);
  
  % Total
  \node[font=\footnotesize, accent] (total) at (9, -3) {$\mathcal{L} = (1-\lambda)\cdot\mathcal{L}_\text{pred} + \lambda\cdot\text{SIGReg}$};
  \draw[->, neutral, thick] (inv.south) |- (total.west);
  \draw[->, neutral, thick] (sig.south) |- (total.east);
\end{tikzpicture}
\vspace{-5pt}
\begin{center}
  \sout{\textcolor{bad}{No stop-gradient / No teacher-student / No EMA / No register tokens}}
\end{center}
\takeaway{\alert{1 loss $\cdot$ 1 hyperparameter $\cdot$ 0 heuristic}}
\end{frame}

% ============================================================
%  SLIDE 20 — LeJEPA vs Prior JEPA (Layout F)
% ============================================================
\begin{frame}{Tại Sao LeJEPA Tốt Hơn DINO / I-JEPA}
\centering
\resizebox{0.95\linewidth}{!}{
\begin{tabular}{p{3.8cm} c c c}
\toprule
\textbf{Tiêu chí} & \textbf{DINO} & \textbf{I-JEPA} & \textbf{LeJEPA} \\
\midrule
Anti-collapse & \cellcolor{bad!20}EMA + stop-grad \textcolor{bad}{\xmark} & \cellcolor{bad!20}Stop-grad + mask \textcolor{bad}{\xmark} & \cellcolor{good!20}\textbf{Theorem} \textcolor{good}{\cmark} \\
\# Hyperparameters & \cellcolor{bad!20}$\geq$7 \textcolor{bad}{\xmark} & \cellcolor{bad!20}$\geq$5 \textcolor{bad}{\xmark} & \cellcolor{good!20}\textbf{1 ($\lambda$)} \textcolor{good}{\cmark} \\
Complexity & \cellcolor{neutral!20}$\mathcal{O}(N)$ & \cellcolor{neutral!20}$\mathcal{O}(N)$ & \cellcolor{good!20}$\mathcal{O}(N)$ \textcolor{good}{\cmark} \\
Architecture support & \cellcolor{bad!20}ViT only \textcolor{bad}{\xmark} & \cellcolor{bad!20}ViT preferred \textcolor{bad}{\xmark} & \cellcolor{good!20}\textbf{Any (60+)} \textcolor{good}{\cmark} \\
Loss informativeness & \cellcolor{bad!20}Low corr. \textcolor{bad}{\xmark} & \cellcolor{bad!20}Low corr. \textcolor{bad}{\xmark} & \cellcolor{good!20}\textbf{94\% Spearman} \textcolor{good}{\cmark} \\
Theory foundation & \cellcolor{bad!20}Post-hoc \textcolor{bad}{\xmark} & \cellcolor{bad!20}Post-hoc \textcolor{bad}{\xmark} & \cellcolor{good!20}\textbf{First-principles} \textcolor{good}{\cmark} \\
Core code (lines) & \cellcolor{bad!20}1000+ \textcolor{bad}{\xmark} & \cellcolor{bad!20}500+ \textcolor{bad}{\xmark} & \cellcolor{good!20}\textbf{$\approx$50} \textcolor{good}{\cmark} \\
\bottomrule
\end{tabular}
}
\vspace{10pt}
\takeaway{\alert{Đơn giản hơn + lý thuyết chắc hơn + work trên mọi kiến trúc}}
\end{frame}

% ============================================================
%  SLIDE 21 — Lambda Trade-off (Layout B)
% ============================================================
\begin{frame}{$\lambda$: Nút Duy Nhất Cần Chỉnh}
\begin{columns}[c]
\column{0.50\textwidth}
  \centering
  \begin{tikzpicture}[scale=0.8, transform shape, every node/.style={font=\scriptsize}]
    % Spectrum line
    \draw[->, thick, neutral] (0,0) -- (6,0) node[right] {$\lambda$};
    
    % Lambda = 0
    \draw[thick] (0.5,0.2) -- (0.5,-0.2) node[below] {$0$};
    \node[bad, align=center, font=\tiny] at (0.5, 1.2) {Pure Invariance\\ \textbf{Collapse}};
    \fill[bad!20] (0.5, 0.5) circle (0.2cm);
    \fill[bad] (0.5, 0.5) circle (0.05cm);
    
    % Lambda = 1
    \draw[thick] (5.5,0.2) -- (5.5,-0.2) node[below] {$1$};
    \node[bad, align=center, font=\tiny] at (5.5, 1.2) {Pure SIGReg\\ \textbf{Uniform}};
    \foreach \px/\py in {5.3/0.6, 5.7/0.4, 5.4/0.3, 5.6/0.7, 5.5/0.5} {
      \fill[bad] (\px,\py) circle (0.03cm);
    }
    
    % Lambda = 0.05
    \draw[thick, good] (2.5,0.2) -- (2.5,-0.2) node[below] {$0.05$};
    \node[good, align=center, font=\tiny] at (2.5, 1.2) {Sweet Spot\\ \textbf{Optimal}};
    \fill[good!20] (2.5, 0.5) circle (0.3cm);
    \foreach \px/\py in {2.4/0.6, 2.6/0.4, 2.3/0.4, 2.7/0.6, 2.5/0.5} {
      \fill[good] (\px,\py) circle (0.04cm);
    }
    
    % Bracket
    \draw[decorate, decoration={brace, amplitude=5pt}, good, thick] 
      (1.5,-0.6) -- (3.5,-0.6) node[midway, below=5pt, font=\tiny, align=center] {Vùng $\lambda$ hoạt động tốt};
  \end{tikzpicture}
\column{0.48\textwidth}
  \centering
  \begin{tikzpicture}
    \begin{axis}[
      width=5.5cm, height=3.5cm,
      xmode=log,
      xlabel={$\lambda$ (log scale)},
      ylabel={Linear Probe Acc (\%)},
      xlabel style={font=\tiny},
      ylabel style={font=\tiny, yshift=-2pt},
      tick label style={font=\tiny},
      xmin=0.005, xmax=0.5,
      ymin=60, ymax=90,
      grid=major, grid style={gray!20},
    ]
      \addplot[color=primary, thick, domain=0.005:0.5, samples=50] {85 - 20*(log10(x) - log10(0.05))^2};
      \addplot[mark=*, good] coordinates {(0.05, 85)} node[above, font=\tiny] {Optimum};
      \fill[good!20, opacity=0.4] (axis cs:0.02,60) rectangle (axis cs:0.1,90);
    \end{axis}
  \end{tikzpicture}
  \vspace{5pt}
  \begin{block}{Recommended defaults}
    \tiny
    $\lambda = 0.05$, $V_g = 2$, $V_l = 6$, batch size $\geq 128$, 
    num\_slices $= 1024$, integration $t \in [-5, 5]$, 17 quadrature points
  \end{block}
\end{columns}
\takeaway{\alert{$\lambda$ robust trên dải rộng — không cần grid search}}
\end{frame}

% ============================================================
%  SLIDE 22 — Informative Loss (Layout B)
% ============================================================
\begin{frame}{Bonus: Training Loss Dự Đoán Được Downstream Performance}
\begin{columns}[T]
\column{0.48\textwidth}
  \vspace{0pt}
  \centering
  \begin{tikzpicture}
    \begin{axis}[
      width=5.0cm, height=4.0cm,
      xlabel={LeJEPA Training Loss},
      ylabel={Linear Probe Acc (\%)},
      xlabel style={font=\tiny},
      ylabel style={font=\tiny, yshift=-2pt},
      tick label style={font=\tiny},
      xmin=0.1, xmax=0.9,
      ymin=60, ymax=90,
      grid=major, grid style={gray!20},
    ]
      \foreach \px/\py in {
        0.2/88, 0.25/86, 0.3/85, 0.35/83, 0.4/81, 0.45/78, 0.5/75, 0.55/73, 0.6/70, 0.7/65, 0.8/62,
        0.22/87, 0.28/84, 0.38/82, 0.42/79, 0.48/76, 0.52/74, 0.65/68, 0.75/63
      } {
        \addplot[only marks, mark=*, mark size=1.5pt, primary!60] coordinates {(\px,\py)};
      }
      \addplot[thick, accent, domain=0.15:0.85] {95 - 40*x};
      \node[font=\tiny, accent, fill=white, inner sep=1pt] at (axis cs:0.6, 85) {$\rho_s \approx 94\%$};
    \end{axis}
  \end{tikzpicture}
  \vspace{-5pt}
  \begin{center}
      \tiny \textit{Mỗi điểm = 1 experiment configuration}
  \end{center}
\column{0.48\textwidth}
  \vspace{0pt}
  \begin{block}{Spearman Correlation}
    \tiny
    \[
    C^{(\alpha)} = \rho_s\!\left(\frac{\text{train\_loss}}{\lambda^\alpha}, \text{test\_acc}\right) 
    \]
    \vspace{-5pt}
    \begin{itemize}
        \item $\alpha=0 \implies \approx 94\%$
        \item $\alpha=0.4 \implies \approx 99\%$
    \end{itemize}
  \end{block}
  \vspace{0pt}
  \begin{alertblock}{Implication}
    \tiny
    Không cần chạy full linear probe để chọn model! Dùng training loss trực tiếp làm proxy $\to$ tiết kiệm \textbf{90\%} compute.
  \end{alertblock}
  \vspace{0pt}
  \scriptsize
  \textbf{So sánh với các phương pháp khác:}
  \begin{itemize}
      \item DINO loss: $\sim 20\%$ correlation
      \item I-JEPA loss: $\sim 30\%$ correlation
      \item \textbf{LeJEPA loss: \hl{94\% correlation}}
  \end{itemize}
\end{columns}
\takeaway{\alert{Lần đầu tiên: SSL có training signal thật sự có nghĩa}}
\end{frame}

% ============================================================
%  SLIDE 23 — Relation to Prior Work (Layout B)
% ============================================================
\begin{frame}{Kết Nối: LeJEPA Generalizes VICReg}
\begin{columns}[T]
\column{0.48\textwidth}
  \vspace{0pt}
  \centering
  \begin{tikzpicture}[
    every node/.style={font=\scriptsize, align=center, rectangle, draw, rounded corners},
    scale=0.9, transform shape
  ]
    \node[draw=primary, fill=primary!10, text=primary, text width=4cm] (root) at (0, 0) {\textbf{LeJEPA}\\(SIGReg + Invariance)};
    
    \node[draw=good, fill=good!10, text=good, text width=4cm] (c1) at (0, -1.8) {\textbf{T = Epps-Pulley}\\ LeJEPA (Recommended)};
    \node[draw=neutral, fill=neutral!10, text=neutral, text width=4cm] (c2) at (0, -3.2) {\textbf{T = mean$^2$ + (std-1)$^2$}\\ VICReg (Special Case)};
    \node[draw=bad, fill=bad!10, text=bad, text width=4cm] (c3) at (0, -4.6) {\textbf{T = Jarque-Bera}\\ Ext. VICReg (Unstable)};
    
    \draw[->, neutral, thick] (root.south) -- (c1.north);
    \draw[->, neutral, thick] (root.west) -- ++(-0.5,0) |- (c2.west);
    \draw[->, neutral, thick] (root.west) -- ++(-0.5,0) |- (c3.west);
  \end{tikzpicture}
\column{0.48\textwidth}
  \vspace{0pt}
  \begin{block}{LeJEPA $\supset$ VICReg}
    \footnotesize
    Khi $T(\mathbf{u}) = \hat\mu(\mathbf{u})^2 + (\hat\sigma(\mathbf{u})^2 - 1)^2$:
    \[
    \text{SIGReg} \xrightarrow{M \to \infty} \text{VICReg variance term}
    \]
    (SIGReg đảm bảo $\mathbb{E}[\mathbf{Z}]=\mathbf{0}$, $\text{Cov}(\mathbf{Z})=I$ trong expectation)
  \end{block}
  \vspace{0pt}
  \begin{alertblock}{Tại sao không dùng VICReg test?}
    \footnotesize
    \textbf{Theorem 4.3:} Finite moments $\implies$ shortcut solutions. \\
    VICReg đã quan sát thấy điều này trong practice (cần các heuristic khác để ổn định).
  \end{alertblock}
\end{columns}
\takeaway{\alert{VICReg là case đặc biệt của LeJEPA — nhưng với test yếu hơn}}
\end{frame}

% ============================================================
%  SLIDE 24 — Stability (Layout B)
% ============================================================
\begin{frame}{LeJEPA: Stable Trên 50+ Architectures Out-of-the-Box}
\begin{columns}[c]
\column{0.58\textwidth}
  \centering
  \begin{tikzpicture}
  \begin{axis}[
    width=7.5cm, height=6.5cm,
    xlabel={Top-1 Accuracy (\%)},
    ylabel={Kiến trúc},
    xlabel style={font=\tiny},
    ylabel style={font=\tiny},
    tick label style={font=\tiny},
    xmin=91.0, xmax=95.5,
    ymin=0, ymax=9,
    ytick={1,2,3,4,5,6,7,8},
    yticklabels={ResNet, ViT, ConvNeXt, Swin, MaxViT, EfficientNet, MobileNet, RegNet},
    legend style={font=\tiny, at={(0.02,0.98)}, anchor=north west},
    grid=major, grid style={gray!20},
    scatter/classes={
      a={mark=*, primary},
      b={mark=*, accent},
      c={mark=*, good},
      d={mark=*, bad},
      e={mark=*, neutral},
      f={mark=*, cyan},
      g={mark=*, magenta},
      h={mark=*, orange}
    }
  ]
    \addplot[scatter, only marks, scatter src=explicit symbolic]
      coordinates {
        (92.5,1) [a] (93.1,1) [a] (93.8,1) [a] (94.2,1) [a]
        (94.5,2) [b] (94.8,2) [b] (95.0,2) [b]
        (93.9,3) [c] (94.4,3) [c] (94.6,3) [c]
        (94.1,4) [d] (94.5,4) [d]
        (94.0,5) [e] (94.6,5) [e]
        (91.5,6) [f] (92.2,6) [f]
        (92.0,7) [g] (92.5,7) [g]
        (93.5,8) [h] (94.1,8) [h]
      };
    \draw[accent, thick, densely dotted] (axis cs:91.5,0) -- (axis cs:91.5,9);
    \draw[accent, thick, densely dotted] (axis cs:95.0,0) -- (axis cs:95.0,9);
    \draw[<->, accent, thick] (axis cs:91.5, 8.5) -- (axis cs:95.0, 8.5) node[midway, above, font=\tiny, text=accent] {Spread 3.5\%};
  \end{axis}
  \end{tikzpicture}
\column{0.40\textwidth}
  \begin{block}{Kết Quả}\footnotesize
    \begin{itemize}
        \item 50 models từ 8 families
        \item Params: 5M $\to$ 700M
        \item Best: 95\% (ViT)
        \item Worst: 91.5\% (EfficientNet)
        \item \textbf{Không có model collapse}
    \end{itemize}
  \end{block}
  \vspace{4pt}
  \begin{exampleblock}{Ablation (Stable)}\footnotesize
    \begin{itemize}
        \item Batch size $128 \to 1024$: $72\% \to 74\%$
        \item Slices $512 \to 4096$: $74\% \to 75\%$
        \item Integration $[-1,1] \to [-5,5]$: nhỏ
    \end{itemize}
  \end{exampleblock}
\end{columns}
\takeaway{50 models, zero collapses — không cần architecture-specific tuning}
\end{frame}

% ============================================================
%  SLIDE 25 — Scale: ViT-H/14 trên ImageNet-1K (Layout H)
% ============================================================
\begin{frame}{LeJEPA Scales: ViT-H/14 Đạt 79\% Trên ImageNet-1K}
\vspace{0.2cm}
\begin{center}
  \begin{tcolorbox}[
    colback=primary!10, colframe=primary,
    width=0.55\textwidth, halign=center, arc=4pt, boxrule=1pt
  ]
    {\Huge\bfseries\color{primary} 79.0\%}\\[4pt]
    {\normalsize ImageNet-1K (linear probe)}\\
    {\footnotesize\color{neutral} ViT-H/14, 100 epochs}
  \end{tcolorbox}
\end{center}
\vspace{0.2cm}
\begin{columns}[c]
\column{0.48\textwidth}
  \centering
  \begin{tikzpicture}
  \begin{axis}[
    width=6.5cm, height=3.5cm,
    ybar,
    bar width=15pt,
    enlarge x limits=0.4,
    ylabel={Top-1 Acc. (\%)},
    symbolic x coords={I-JEPA ViT-H, LeJEPA ViT-L, LeJEPA ConvNeXt-H},
    xtick=data,
    nodes near coords,
    nodes near coords align={vertical},
    every node near coord/.append style={font=\tiny},
    ymin=75, ymax=80,
    xticklabel style={font=\tiny, text width=1.5cm, align=center},
    ylabel style={font=\tiny},
    tick label style={font=\tiny},
    grid=major, grid style={gray!20}
  ]
    \addplot[fill=neutral!60, draw=none] coordinates {(I-JEPA ViT-H,77.3)};
    \addplot[fill=primary!80, draw=none] coordinates {(LeJEPA ViT-L,78.2)};
    \addplot[fill=good!80, draw=none] coordinates {(LeJEPA ConvNeXt-H,79.0)};
  \end{axis}
  \end{tikzpicture}
\column{0.48\textwidth}
  \centering\footnotesize
  \textbf{Transfer Learning}\\
  \vspace{4pt}
  \begin{tabular}{p{2.3cm} r r r r}
  \toprule
  & \multicolumn{2}{c}{1-shot} & \multicolumn{2}{c}{10-shot} \\
  Method & DTD & Avg & DTD & Avg \\
  \midrule
  I-JEPA ViT-H\newline(300ep) & 27.7 & 30.2 & 57.7 & 60.5 \\
  \textbf{LeJEPA ViT-L\newline(100ep)} & \textbf{33.2} & \textbf{29.6} & \textbf{64.7} & \textbf{61.0} \\
  \bottomrule
  \end{tabular}
\end{columns}
\takeaway{3x ít compute, 2x nhỏ hơn — vẫn SOTA}
\end{frame}

% ============================================================
%  SLIDE 26 — In-Domain: Đánh bại DINOv2 trên Galaxy10
% ============================================================
\begin{frame}{In-Domain SSL Đánh Bại Frontier Models}
\begin{columns}[c]
\column{0.48\textwidth}
  \centering\footnotesize\textbf{Frozen Backbone}
  \vspace{2pt}
  
  \begin{tikzpicture}
  \begin{axis}[
    width=6.5cm, height=4.5cm,
    xlabel={Shots per class},
    ylabel={Accuracy (\%)},
    xlabel style={font=\tiny},
    ylabel style={font=\tiny},
    tick label style={font=\tiny},
    xmin=0, xmax=50,
    ymin=40, ymax=90,
    legend style={font=\tiny, at={(0.98,0.05)}, anchor=south east, draw=none, fill=lightgray, inner sep=2pt},
    grid=major, grid style={gray!20}
  ]
    \addplot[color=neutral!50, thick, mark=square] coordinates {(1,45) (5,55) (10,65) (50,75)};
    \addlegendentry{DINOv2 ViT-S}
    \addplot[color=neutral!80, thick, mark=triangle] coordinates {(1,48) (5,58) (10,68) (50,78)};
    \addlegendentry{DINOv3 ViT-S}
    \addplot[color=good, very thick, mark=*] coordinates {(1,55) (5,68) (10,75) (50,85)};
    \addlegendentry{LeJEPA ResNet-34 \cmark}
  \end{axis}
  \end{tikzpicture}
\column{0.48\textwidth}
  \centering\footnotesize\textbf{Full Finetuning}
  \vspace{2pt}
  
  \begin{tikzpicture}
  \begin{axis}[
    width=6.5cm, height=4.5cm,
    xlabel={Shots per class},
    ylabel={Accuracy (\%)},
    xlabel style={font=\tiny},
    ylabel style={font=\tiny},
    tick label style={font=\tiny},
    xmin=0, xmax=50,
    ymin=50, ymax=95,
    legend style={font=\tiny, at={(0.98,0.05)}, anchor=south east, draw=none, fill=lightgray, inner sep=2pt},
    grid=major, grid style={gray!20}
  ]
    \addplot[color=neutral!50, thick, mark=square] coordinates {(1,55) (5,65) (10,75) (50,85)};
    \addlegendentry{DINOv2 ViT-S}
    \addplot[color=neutral!80, thick, mark=triangle] coordinates {(1,58) (5,68) (10,78) (50,88)};
    \addlegendentry{DINOv3 ViT-S}
    \addplot[color=good, very thick, mark=*] coordinates {(1,65) (5,75) (10,85) (50,92)};
    \addlegendentry{LeJEPA ResNet-34 \cmark}
  \end{axis}
  \end{tikzpicture}
\end{columns}
\vspace{0.2cm}
\begin{exampleblock}{Paradigm Shift}\footnotesize
  DINOv2/v3 train trên hàng triệu ảnh tự nhiên. Galaxy10 là ảnh thiên hà (domain shift lớn).\\
  Không cần chờ frontier model cho domain của bạn. Train LeJEPA trực tiếp trên data domain-specific với 1 GPU.
\end{exampleblock}
\takeaway{Domain-specific SSL > Generic transfer — ngay cả với dataset 11K samples}
\end{frame}

% ============================================================
%  SLIDE 27 — Emergent Semantic Structure (Layout B)
% ============================================================
\begin{frame}{Emergent Semantics: LeJEPA Học Segment Vật Thể}
\begin{columns}[c]
\column{0.55\textwidth}
  \centering
  \begin{tikzpicture}
    % Original image placeholder
    \node[draw=neutral, fill=neutral!20, minimum width=2.5cm, minimum height=3cm, font=\small, align=center] (img) at (0,0) {Ảnh gốc\\(Chó trên cỏ)};
    
    % Arrow
    \draw[->, thick, primary] (img.east) -- ++(1,0) node[midway, above, font=\tiny] {PCA};
    
    % PCA features placeholder with shading
    \draw[draw=primary, thick] (2.3,-1.5) rectangle (4.8,1.5);
    % Background shading (cool colors)
    \fill[left color=blue!30, right color=gray!30] (2.3,-1.5) rectangle (4.8,1.5);
    % Foreground object shading (warm colors)
    \fill[left color=red!60, right color=orange!60] (3.55,-0.5) ellipse (0.8cm and 1cm);
    
    \node[font=\small, text=white] at (3.55,-0.5) {Vật thể};
    \node[font=\tiny, align=center] at (4.3, 1.0) {Nền};
    
    \node[font=\footnotesize, primary, align=center, yshift=-0.4cm] at (3.55,-1.5) {PCA Features\\(Không supervision)};
  \end{tikzpicture}
\column{0.43\textwidth}
  \begin{block}{Kết Quả Phân Tích PCA}\footnotesize
    \begin{itemize}
        \item Cắt 3 thành phần PCA chính từ layer cuối, ánh xạ sang RGB.
        \item \textbf{Màu ấm (đỏ/cam):}\\Foreground objects
        \item \textbf{Màu lạnh (xanh/xám):}\\Background/foliage
    \end{itemize}
  \end{block}
  \vspace{4pt}
  \begin{alertblock}{Ý Nghĩa}\footnotesize
    Zero labels — LeJEPA tự phân tách cấu trúc ngữ nghĩa qua attention thresholding.
  \end{alertblock}
\end{columns}
\takeaway{SSL objective đúng $\to$ semantic structure tự nổi lên}
\end{frame}

% ============================================================
%  SLIDE 28 — Impact: Bản đồ ảnh hưởng (Layout G / Radial Map)
% ============================================================
\begin{frame}{Impact: LeJEPA Mở Ra Những Gì?}
\vspace{-0.2cm}
\begin{center}
  \begin{tikzpicture}[
      scale=0.75, transform shape,
      center node/.style={circle, draw=primary, fill=primary!80, text=white, font=\bfseries\large, minimum size=2.5cm, align=center},
      branch node/.style={rectangle, draw, rounded corners=4pt, font=\footnotesize, align=center, minimum height=0.8cm, minimum width=2cm},
      timeline node/.style={font=\scriptsize, align=center}
    ]
    % Center Node
    \node[center node] (center) at (0, 0) {LeJEPA};
    
    % 5 Branches using \foreach
    \foreach \angleNode/\lbl/\icon/\clr in {
      90/Robotics\\World Models/Robot/accent,
      20/Scientific\\ML/Telescope/good,
      -30/Edge\\AI/Chip/primary,
      210/Open\\Ecosystem/GitHub/neutral,
      150/Academic\\Theory/Grad Cap/bad%
    } {
      \node[branch node, draw=\clr, fill=\clr!10, text=\clr] (n\angleNode) at (\angleNode:3.8cm) {\textbf{\icon}\\[2pt]\lbl};
      \draw[->, thick, \clr] (center) -- (n\angleNode);
    }
  \end{tikzpicture}
\end{center}

\vspace{0.3cm}
\begin{center}
  \begin{tikzpicture}[scale=0.8, transform shape, every node/.style={font=\scriptsize, align=center}]
    % Timeline
    \draw[->, thick, neutral] (0,0) -- (10,0);
    
    \foreach \posX/\yearL/\method/\desc in {
      0/2020/SimCLR/(heuristic),
      2.5/2021/DINO/(scale),
      5/2023/DINOv2/(scale$^2$),
      7.5/2024/I-JEPA/(masking),
      10/2025/LeJEPA $\star$/(theory)%
    } {
      \draw[thick, neutral] (\posX, 0.1) -- (\posX, -0.1);
      \node[above] at (\posX, 0.1) {\textbf{\yearL}};
      \node[below] at (\posX, -0.1) {\textbf{\method}\\[2pt]\desc};
    }
  \end{tikzpicture}
\end{center}

\vspace{-0.3cm}
\takeaway{\alert{Nếu 2024 là năm của scaling laws, 2025 là năm tái khám phá structure}}
\end{frame}

% ============================================================
%  SLIDE 29 — Sử dụng thực tế & Citation (Layout B)
% ============================================================
\begin{frame}[fragile]{Bắt Đầu Dùng LeJEPA}
\begin{columns}[c]
\column{0.48\textwidth}
  \begin{center}
    \begin{tikzpicture}[
      node distance=0.8cm,
      mynode/.style={rectangle, draw=primary, fill=primary!10, rounded corners=4pt, font=\scriptsize, align=center, minimum width=3.5cm, minimum height=0.6cm},
      center node/.style={circle, draw=primary, fill=primary!80, text=white, font=\bfseries\small, minimum size=1.5cm, align=center}
    ]
      \node[center node] (center) at (0, 0) {LeJEPA};
      
      \node[mynode, left=1.2cm of center] (github) {\textbf{GitHub}\\rbalestr-lab/lejepa\\50 LOC core};
      \node[mynode, above=1cm of center] (integ) {\textbf{Integration}\\timm + Lightning};
      \node[mynode, right=1.2cm of center] (eval) {\textbf{Evaluation}\\OnlineProbe, RankMe};
      \node[mynode, below=1cm of center] (cite) {\textbf{Citation}\\arXiv: 2511.08544};
      
      \draw[->, thick, primary] (center) -- (github);
      \draw[->, thick, primary] (center) -- (integ);
      \draw[->, thick, primary] (center) -- (eval);
      \draw[->, thick, primary] (center) -- (cite);
    \end{tikzpicture}
  \end{center}
\column{0.48\textwidth}
  \begin{block}{Quick Start}
\begin{lstlisting}
import lejepa

# Initialize Epps-Pulley test
test = lejepa.multivariate.SlicingUnivariateTest(
    univariate_test=lejepa.univariate.EppsPulley(17),
    num_slices=1024
)
# Compute loss
sigreg_loss = test(embeddings)
\end{lstlisting}
  \end{block}
  \vspace{5pt}
  \begin{alertblock}{Citation}
    \tiny
    @misc\{balestriero2025lejepa,\\
      \ \ title=\{LeJEPA: Provable and Scalable Self-Supervised Learning\},\\
      \ \ author=\{Balestriero, Randall and LeCun, Yann\},\\
      \ \ year=\{2025\}, eprint=\{2511.08544\}\\
    \}
  \end{alertblock}
\end{columns}
\takeaway{\alert{\texttt{pip install lejepa} — bắt đầu trong 5 phút}}
\end{frame}

% ============================================================
%  SLIDE 30 — Kết luận (Layout A - Hero)
% ============================================================
\begin{frame}[plain]
\vspace{0.5cm}
\begin{center}
  {\Huge\bfseries\color{primary} Kết Luận}\\[16pt]
  
  \begin{tikzpicture}[
      scale=0.85, transform shape,
      step node/.style={rectangle, draw=primary, fill=primary!10, rounded corners=6pt, font=\scriptsize, align=center, minimum width=2.8cm, minimum height=1.2cm},
      final node/.style={rectangle, draw=good, fill=good!20, rounded corners=6pt, font=\footnotesize\bfseries, align=center, minimum width=3.2cm, minimum height=1.4cm}
    ]
    \node[step node] (s1) at (0, 0) {Supervised /\\Heuristic SSL\\[2pt]\color{bad}\textit{Fragile, expensive}};
    \node[step node] (s2) at (4.5, 0) {JEPA (2022)\\[2pt]\color{neutral}\textit{No theory}};
    \node[final node] (s3) at (9.5, 0) {\color{good}LeJEPA (2025)\\[2pt]\color{primary}Theory + Practice};
    
    \draw[->, thick, primary, >=stealth] (s1) -- (s2);
    \draw[->, thick, primary, >=stealth] (s2) -- (s3);
  \end{tikzpicture}
\end{center}

\vspace{0.4cm}
\begin{columns}[c]
\column{0.15\textwidth}
\column{0.70\textwidth}
  \begin{itemize}
    \item[\textcolor{good}{\cmark}] \textbf{THEOREM:} Iso Gaussian uniquely minimizes downstream risk
    \item[\textcolor{good}{\cmark}] \textbf{METHOD:} SIGReg = sliced Epps-Pulley, $\mathcal{O}(N)$, bounded gradient
    \item[\textcolor{good}{\cmark}] \textbf{RESULT:} SOTA trên 60+ archs, beats DINOv2 in-domain
  \end{itemize}
\column{0.15\textwidth}
\end{columns}

\vspace{0.4cm}
\begin{center}
\Large\itshape\color{accent}
``If 2024 was the year of scaling laws,\\
2025 may be the year we rediscover structure.''
\end{center}

\vspace{0.4cm}
\begin{center}
  {\Large\bfseries\color{primary} Q\&A}
\end{center}
\end{frame}

\end{document}
